% Cote 119: offprints, reports, courses, CV and letters (typed and printed part): archive cote 119, site page 67 (PDF page 68).
% offprint. Draft compared: none.
% AI-transcribed by Claude Opus 5.5 (claude-opus-5-5, adaptive thinking, effort high), Message Batches API, batch r119_full (msgbatch_01DwZNm5Y7YsKBWyAN4sFs8w), 2026-11-13; not yet checked by eye.
\scanpage{119/67}
\mspage{( 3 )}

Exemple 2.3. --- Soit $(G, w, t)$ un triple de Tate neutralisé (voir 1.4),
et soit $C \in G_0(\mathbf{R})$ vérifiant
\[ C^2 = \varepsilon = w(-1). \]

Si V est un G-module homogène de poids $n$ une forme $\varphi : V \otimes V \to T^{\otimes -n}$
est appelée de C-polarisation si la forme bilinéaire $\varphi_C$ sur le $\mathbf{R}$-vectoriel V
\[ \varphi_C(x, y) = \varphi(x, C\,y) \]
est symétrique définie positive. Si pour chaque V il existe une forme
de C-polarisation, l'ensemble de ces formes définit une polarisation $\pi_C$
du triple $(G, w, t)$ ; ont dit alors que C est un élément hodgien de $G(\mathbf{R})$
et les polarisations obtenues ainsi sont appelées hodgiennes. Le triple
$(G, w, t)$ est dit hodgien s'il existe $C \in G(\mathbf{R})$ hodgien. On déduit faci-%
lement de 2.2 et de [(${}^1$), 2.5] que C est hodgien si et seulement si le
$\mathbf{R}$-groupe $(G_0)_C$ (voir loc. cit.) est compact et que $(G, w, t)$ est hodgien
si et seulement si $G_0$ est réductif, est une forme tordue intérieure de sa
forme compacte, et si l'invariant $\varepsilon_{C_0}$ ($C_0 = Rep_0(G_0)$, voir [(${}^1$), 3.2])
est égal à $\varepsilon$ modulo $Z_0(\mathbf{R})^2$.

On déduit de 1.3, 2.2 et de (${}^1$) le résultat suivant, qui résume la
théorie des polarisations dans le cas gradué :

Théorème 2.4. --- Soient L un lien algébrique affine sur $\mathbf{R}$ de centre Z,
$t : L \to \mathbf{G}_m$, $w : \mathbf{G}_m \to Z$ des morphismes vérifiant $t \circ w = -2$. Alors :

(a) Si $(C, w, T)$ est un triple de Tate lié par $(L, w, T)$, Pol $(C, w, T)$
est un pseudo-torseur sous ${}_2Z_0(\mathbf{R})$ ;

(b) Si $(C, w, T)$, $(C', w, T')$ sont des triples de Tate liés par $(L, w, t)$
munis de polarisations $\pi$, $\pi'$ il existe une équivalence $(C, w, T) \simeq (C', w, T')$
unique à isomorphisme (non unique) près respectant les polarisations données.

(c) Supposons que L soit connexe ou abélien. Alors $(L, w, t)$ est pola-%
risable [i. e. il existe $(C, w, T)$ lié par $(L, w, t)$ avec Pol $(C, w, T) \neq \emptyset$]
si et seulement si $L_0 = \operatorname{Ker}(L \to G_m)$ (on peut prouver l'existence de
ce noyau) est le lien d'un $\mathbf{R}$-groupe compact. S'il en est ainsi, le triple
$(C, w, T)$ est neutre si et seulement si $\varepsilon = w(-1)$ est un carré dans $Z_0(\mathbf{R})$
et, dans ce cas, toute polarisation de $(C, w, T)$ est hodgienne.

(*) Séance du 8 décembre 1971.
(${}^1$) N. Saavedra, Comptes rendus, 273, série A, 1971, p. 1114.
(${}^2$) P. Deligne. Travaux de Griffiths, Séminaire Bourbaki, 376 (mai-juin 1970).
(${}^3$) Travail en cours de préparation sur les catégories tannakiennes (à paraître dans
les Lectures Notes in Mathematics, Springer-Verlag).
(${}^4$) N. Saavedra, Comptes rendus, 272, série A, 1971, p. 389.

Institut
des Hautes Études Scientifiques,
35, route de Chartres,
91-Bures-sur-Yvette,
Essonne.

\mspage{\hand{67}}
